\section{Limitations and discussion}
\label{sec:limitations}

\paragraph{Single-layer agents.} Each agent is a perceptron. This is a deliberate
choice --- it is the simplest model in which two agents can agree on some tasks and
disagree on others, and it keeps the update rule in closed form, which is what
makes the mechanism identifiable at all. It is nonetheless the obvious objection,
and \S\ref{sec:phase} carries one experiment outside it.
\todo{make sure that experiment exists before submission; without it this
paragraph is an admission rather than a limitation.}

\paragraph{The perturbation is injected, not learned.} We shift the opinion field
by a constant and study the consequences. We do not model how an agent would come
to encode messages in a source-dependent way in the first place. That is a
separate question and we make no claim about it.

\paragraph{Scale.} $N = 40$ agents, and the balance observables cost $O(N^3)$
before the closed form of \S\ref{sec:observables}. The finite-size scaling in
\S\ref{sec:boundary} is what licenses reading the boundary as a genuine transition
rather than a feature of a small population, and it is the analysis to trust over
any single map.

\paragraph{The dynamics anneals rather than equilibrating.} $F_C, F_V < 0$ almost
everywhere, so the population slows down instead of reaching a stationary state.
Every observable is therefore reported at a stated measurement time, and the
qualitative structure --- which regime, which side of the diagonal --- is stable
across the range we checked, while the numerical values are not
time-independent. We regard time-to-organise as part of the phenomenon rather
than a nuisance; \S\ref{sec:agenda} is entirely a statement about it.

\paragraph{On the social reading.} These are single-layer networks, and extending
any of this to human behaviour is not automatic. We use the words group, trust and
discrimination because they name the structures the model actually contains --- a
label carrying no task information, a learned per-source reliability, and a
partition aligned with the label --- and not as a claim about people. What the
result does support is narrower and, we think, more useful: a mechanism by which
systems that learn from each other can acquire and entrench a partition along a
feature that no component of the system was ever told to care about, without
biased data and without a preference anywhere in the system.

\paragraph{Objective misspecification by scale.} The update rule is optimal for a
teacher--student pair. The two properties that make it good at that job --- moving
away from sources it has learned to distrust, and attributing surprise to whichever
sector is less certain --- are exactly the two that produce the collective
pathology. This is not a defect in the derivation. It is what happens when an
objective derived at one scale is deployed at another, and it suggests that the
right question to ask of a peer-weighting scheme is not whether it is optimal for a
pair but what it does to a population. Ablation A2 gives a concrete instance:
declining to anti-learn from unreliable peers costs sample efficiency and buys
immunity to this failure mode.
